\section{Images and Inverse Images}
\chapterline

Up to this point, we have watched functions act on individual points: you feed an input $x$ into the machine, and an output $f(x)$ emerges on the other side. 

In higher mathematics, however, functions rarely operate one particle at a time. They act on entire regions, subspaces, and collections. A function $f: X \to Y$ creates a two-way transportation network between the subsets of $X$ and the subsets of $Y$:
\begin{itemize}
\item \textbf{Pushing Forward (Image):} Take a subset $S \subseteq X$ and transport it forward into $Y$ by collecting all outputs $f(s)$.
\item \textbf{Pulling Backward (Inverse Image / Preimage):} Select a target region $T \subseteq Y$ and trace backward to collect all points in $X$ whose outputs land in $T$.
\end{itemize}

At first glance, you might expect pushing forward and pulling backward to be mirror images of each other. But a deep asymmetry lives here: \emph{pulling backward behaves dramatically better than pushing forward}. Preimages preserve unions, intersections, and complements perfectly. Forward images, by contrast, can fracture and lose information whenever a function merges distinct inputs.

\begin{historybox}
This profound algebraic asymmetry was uncovered in the early twentieth century by Felix Hausdorff (\emph{Grundz\"uge der Mengenlehre}, 1914) alongside pioneers of the Polish and French schools of analysis (Maurice Fr\'echet, Kazimierz Kuratowski, and Henri Lebesgue). 

Mathematicians were trying to liberate calculus from the rigid confines of Euclidean $\R^n$. In abstract spaces where there are no coordinates, no rulers, and no $\eps$-$\delta$ formulas, how can one define whether a transformation is ``continuous'' or whether a region is ``measurable''? 

The early naive instinct was to push sets forward: if $U$ is an open neighborhood, shouldn't its image $f(U)$ also be open? But that approach collapsed immediately into counterexamples. Hausdorff realized that the true invariant structure of mathematics flows backward: the inverse image operation $f^{-1}$ is an exact Boolean algebra homomorphism. It preserves arbitrary unions, intersections, and complements without exception. This single algebraic discovery crowned the preimage as the primary structural engine for modern topology, measure theory, and modern analysis.
\end{historybox}

\subsection{Formal Definitions}

\begin{definition}[Forward Image]
Let $f: X \to Y$ be a function and let $S \subseteq X$. The \textbf{image} of $S$ under $f$ is
\[
f(S) = \{f(s) \in Y : s \in S\} = \{y \in Y : (\exists s \in S)(f(s) = y)\}.
\]
\end{definition}

\begin{definition}[Inverse Image / Preimage]
Let $f: X \to Y$ be a function and let $T \subseteq Y$. The \textbf{inverse image} (or \textbf{preimage}) of $T$ under $f$ is
\[
f^{-1}(T) = \{x \in X : f(x) \in T\}.
\]
\end{definition}

\begin{center}
\begin{tikzpicture}[scale=0.92, >=Latex]
  % Panel 1: Forward Push with Information Loss
  \begin{scope}[yshift=0cm]
    \node[anchor=west, font=\small\bfseries] at (-0.5, 2.75) {Pushing Forward: $f(S_1 \cap S_2) \subsetneq f(S_1) \cap f(S_2)$ (Collisions Merge)};
    \draw[rounded corners=4pt, fill=gray!3, draw=gray!25] (-0.6, -0.6) rectangle (13.4, 2.5);
    
    % Domain X
    \draw[thick, sourceblue, fill=sourceblue!8] (1.5, 0.9) ellipse (1.4cm and 1.1cm);
    \node[sourceblue!90!black, font=\scriptsize\bfseries] at (1.5, 2.15) {Domain $X$};
    
    % Disjoint sets S1 and S2
    \draw[thick, trapred!70, fill=trapred!15] (0.9, 0.9) circle (0.42cm);
    \node[trapred!90!black, font=\scriptsize\bfseries] at (0.9, 0.9) {$S_1$};
    \draw[thick, sourceblue!70, fill=sourceblue!20] (2.1, 0.9) circle (0.42cm);
    \node[sourceblue!90!black, font=\scriptsize\bfseries] at (2.1, 0.9) {$S_2$};
    
    % Codomain Y
    \draw[thick, targetgreen, fill=targetgreen!8] (11.2, 0.9) ellipse (1.6cm and 1.1cm);
    \node[targetgreen!90!black, font=\scriptsize\bfseries] at (11.2, 2.15) {Codomain $Y$};
    
    % Overlapping images
    \draw[thick, trapred!70, fill=trapred!15] (10.8, 0.9) circle (0.52cm);
    \draw[thick, sourceblue!70, fill=sourceblue!20] (11.6, 0.9) circle (0.52cm);
    \node[trapred!90!black, font=\tiny\bfseries] at (10.45, 0.9) {$f(S_1)$};
    \node[sourceblue!90!black, font=\tiny\bfseries] at (11.95, 0.9) {$f(S_2)$};
    
    % Arrows
    \draw[->, thick, trapred!80] (1.35, 1.25) to[bend left=16] node[above, font=\scriptsize] {$f(S_1)$} (10.25, 1.25);
    \draw[->, thick, sourceblue!80] (2.55, 0.55) to[bend right=16] node[below, font=\scriptsize] {$f(S_2)$} (11.05, 0.45);
    
    \node[anchor=center, font=\scriptsize, align=center] at (6.3, 0.9) {
      $S_1 \cap S_2 = \emptyset \implies f(S_1 \cap S_2) = \emptyset$\\
      Yet $f(S_1) \cap f(S_2) \neq \emptyset$!\\
      Collisions create false overlap.
    };
  \end{scope}

  % Panel 2: Pulling Backward Preserves Structure
  \begin{scope}[yshift=-3.9cm]
    \node[anchor=west, font=\small\bfseries] at (-0.5, 2.75) {Pulling Backward: $f^{-1}(T_1 \cap T_2) = f^{-1}(T_1) \cap f^{-1}(T_2)$ (Exact Match)};
    \draw[rounded corners=4pt, fill=targetgreen!3, draw=targetgreen!25] (-0.6, -0.6) rectangle (13.4, 2.5);
    
    % Domain X
    \draw[thick, sourceblue, fill=sourceblue!8] (1.5, 0.9) ellipse (1.4cm and 1.1cm);
    \node[sourceblue!90!black, font=\scriptsize\bfseries] at (1.5, 2.15) {Domain $X$};
    
    \draw[thick, quotientorange, fill=quotientorange!20] (1.5, 0.9) circle (0.65cm);
    \node[quotientorange!90!black, font=\tiny\bfseries] at (1.5, 0.9) {$f^{-1}(T_1 \cap T_2)$};
    
    % Codomain Y
    \draw[thick, targetgreen, fill=targetgreen!8] (11.2, 0.9) ellipse (1.6cm and 1.1cm);
    \node[targetgreen!90!black, font=\scriptsize\bfseries] at (11.2, 2.15) {Codomain $Y$};
    
    % Overlapping T1 and T2
    \draw[thick, structurepurple!70, fill=structurepurple!15] (10.8, 0.9) circle (0.52cm);
    \draw[thick, targetgreen!70, fill=targetgreen!20] (11.6, 0.9) circle (0.52cm);
    \node[structurepurple!90!black, font=\tiny\bfseries] at (10.45, 0.9) {$T_1$};
    \node[targetgreen!90!black, font=\tiny\bfseries] at (11.95, 0.9) {$T_2$};
    
    % Reverse arrows
    \draw[<-, thick, structurepurple!80] (2.2, 1.2) to[bend left=16] node[above, font=\scriptsize] {$f^{-1}(T_1)$} (10.25, 1.2);
    \draw[<-, thick, targetgreen!80] (2.2, 0.6) to[bend right=16] node[below, font=\scriptsize] {$f^{-1}(T_2)$} (11.05, 0.5);
    
    \node[anchor=center, font=\scriptsize, align=center] at (6.3, 0.9) {
      $x \in f^{-1}(T_1 \cap T_2) \iff f(x) \in T_1 \cap T_2$\\
      Because every input has a unique output,\\
      logic flows reversibly.
    };
  \end{scope}
\end{tikzpicture}
\end{center}

\begin{trapbox}
The symbol $f^{-1}(T)$ represents a set-theoretic operation on subsets of the codomain. It exists for \textbf{every} function $f$, whether or not $f$ is invertible! If $f(x) = x^2$, the inverse function $f^{-1}: \R \to \R$ does not exist, yet $f^{-1}([1, 4]) = [-2, -1] \cup [1, 2]$ is a perfectly well-defined subset of $\R$.
\end{trapbox}

\subsection{Why Preimages Behave So Beautifully}

To understand why preimages preserve set operations without exception, notice that determining whether $x \in f^{-1}(T)$ comes down to evaluating a single output: does $f(x)$ land inside $T$?

Because $f$ is a function, every element $x$ has exactly one output $f(x)$. There is no ambiguity. The point $f(x)$ either belongs to an intersection, union, or complement, or it does not.

\begin{theorem}[Preimages Distribute Over Intersections]
Let $f: X \to Y$ be a function and let $T_1, T_2 \subseteq Y$. Then
\[
f^{-1}(T_1 \cap T_2) = f^{-1}(T_1) \cap f^{-1}(T_2).
\]
\end{theorem}

\begin{proof}
Let $x \in X$ be arbitrary. We trace the logical biconditionals:
\begin{align*}
x \in f^{-1}(T_1 \cap T_2) 
&\iff f(x) \in T_1 \cap T_2 && \text{(definition of preimage)} \\
&\iff f(x) \in T_1 \text{ and } f(x) \in T_2 && \text{(definition of intersection)} \\
&\iff x \in f^{-1}(T_1) \text{ and } x \in f^{-1}(T_2) && \text{(definition of preimage)} \\
&\iff x \in f^{-1}(T_1) \cap f^{-1}(T_2) && \text{(definition of intersection)}.
\end{align*}
Since each step is an equivalence, the two sets contain identical elements:
\[
f^{-1}(T_1 \cap T_2) = f^{-1}(T_1) \cap f^{-1}(T_2). \qedhere
\]
\end{proof}

The identical argument works for unions and complements:
\[
f^{-1}(T_1 \cup T_2) = f^{-1}(T_1) \cup f^{-1}(T_2), \qquad f^{-1}(Y \setminus T) = X \setminus f^{-1}(T).
\]
The preimage is an algebraic homomorphism on Boolean algebras: it commutes with every standard set operation.

\subsection{Where Forward Images Break Down}

Now contrast this with the forward image of an intersection:

\begin{theorem}
Let $f: X \to Y$ and let $S_1, S_2 \subseteq X$. Then
\[
f(S_1 \cap S_2) \subseteq f(S_1) \cap f(S_2).
\]
\end{theorem}

\begin{proof}
Let $y \in f(S_1 \cap S_2)$ be arbitrary. By definition of the forward image, there exists some $x \in S_1 \cap S_2$ such that $f(x) = y$.

Since $x \in S_1 \cap S_2$, we have $x \in S_1$ and $x \in S_2$.
\begin{itemize}
\item Because $x \in S_1$ and $f(x) = y$, we have $y \in f(S_1)$.
\item Because $x \in S_2$ and $f(x) = y$, we have $y \in f(S_2)$.
\end{itemize}
Thus $y \in f(S_1) \cap f(S_2)$. Hence $f(S_1 \cap S_2) \subseteq f(S_1) \cap f(S_2)$.
\end{proof}

\emph{Why does the reverse inclusion fail?} 
Suppose $y \in f(S_1) \cap f(S_2)$. This means $y \in f(S_1)$ and $y \in f(S_2)$.
Therefore, there exists some $x_1 \in S_1$ with $f(x_1) = y$, and there exists some $x_2 \in S_2$ with $f(x_2) = y$.
\emph{Does this guarantee $x_1 = x_2$?}

No! If $f$ is not injective, $x_1$ and $x_2$ could be two completely different inputs that happen to collide on the same destination $y$. If $x_1 \neq x_2$, neither element needs to live in the intersection $S_1 \cap S_2$, leaving $f(S_1 \cap S_2)$ empty even though $f(S_1) \cap f(S_2)$ is non-empty!

\begin{example}[A Clean Counterexample]
Let $f: \R \to \R$ be $f(x) = x^2$.
Take the disjoint subsets $S_1 = \{-2\}$ and $S_2 = \{2\}$.
Their intersection is empty: $S_1 \cap S_2 = \emptyset$, so
\[
f(S_1 \cap S_2) = f(\emptyset) = \emptyset.
\]
However, computing the images separately:
\[
f(S_1) = \{4\}, \qquad f(S_2) = \{4\}, \implies f(S_1) \cap f(S_2) = \{4\}.
\]
Here $\emptyset \subsetneq \{4\}$. The equality fails precisely because $-2$ and $2$ collided under the squaring map.
\end{example}

\begin{theorem}
The equality $f(S_1 \cap S_2) = f(S_1) \cap f(S_2)$ holds for all subsets $S_1, S_2 \subseteq X$ if and only if $f$ is injective.
\end{theorem}

\subsection{Why Mathematicians Care: The Gateway to Topology}

This structural superiority of the preimage is the reason modern topology defines \emph{continuous functions} via inverse images:
\begin{center}
\emph{A function $f: X \to Y$ is continuous if for every open set $V \subseteq Y$, the preimage $f^{-1}(V)$ is open in $X$.}
\end{center}
If topologists had tried to define continuity using forward images (requiring $f(U)$ to be open whenever $U$ is open), even standard, smoothly behaving functions like $f(x) = x^2$ would fail: the open interval $(-1, 1)$ maps to $[0, 1)$, which includes its boundary point $0$ and is not open in $\R$! Preimages preserve topological structure without distortion, while forward images routinely fold and collapse boundary boundaries.

\subsection{What Preimages Unlocked, Their Limits, and Where We Go Next}

Mastering set transport unlocked the architecture of modern pure mathematics:
\begin{itemize}
\item \textbf{Topology without coordinates:} Defining continuity via $f^{-1}(V)$ (Chapter~22) allows us to study curved spaces, infinite-dimensional function spaces, and manifolds without ever drawing an $x$-$y$ axis.
\item \textbf{Measure theory and probability:} In rigorous probability, a random variable $X$ is not a mystery number; it is a measurable function, and the probability of an event $E$ is evaluated by pulling $E$ backward into the sample space: $\P(X \in E) = \P(X^{-1}(E))$.
\item \textbf{Algebraic pullbacks:} Dual spaces, differential forms, and bundle projections all rely on preimages to pull structures back from target spaces into source spaces.
\end{itemize}

Yet image and preimage operations have an inherent boundary: they are \textbf{transport mechanisms between two pre-existing spaces} $X$ and $Y$. They take sets that already exist in one world and ship them to another. 

What they cannot do is manufacture a \emph{new} mathematical universe out of an existing one. They cannot take a single space $X$ and simplify it from within—identifying points that are redundant, collapsing symmetries, or gluing opposite edges of a square together to create a cylinder or a donut.

To forge new spaces by collapsing redundant information, we need an entirely different kind of mathematical engine. In Chapter~15, we introduce \textbf{Relations and Equivalence}, discovering how the three axioms of equivalence allow us to quotient spaces and build new mathematical objects from the ground up.

\subsection{Exercises}

\begin{enumerate}
\item Let $f: X \to Y$ and let $S \subseteq X$ and $T \subseteq Y$.
\begin{enumerate}[label=(\alph*)]
\item Prove that $S \subseteq f^{-1}(f(S))$. Give a concrete example where equality fails.
\item Prove that $f(f^{-1}(T)) \subseteq T$. Give a concrete example where equality fails.
\end{enumerate}
\item Prove that $f^{-1}(T_1 \cup T_2) = f^{-1}(T_1) \cup f^{-1}(T_2)$ for any $T_1, T_2 \subseteq Y$.
\item Prove that $f(S_1 \cup S_2) = f(S_1) \cup f(S_2)$ for any $S_1, S_2 \subseteq X$. Why does union preserve equality for forward images while intersection does not?
\item Let $f: \R \to \R$ be defined by $f(x) = x^2 - 2x$. Find:
\begin{enumerate}[label=(\alph*)]
\item $f([0, 3])$.
\item $f^{-1}([-1, 3])$.
\item $f^{-1}(\{-5\})$.
\end{enumerate}
\end{enumerate}

\newpage
